\begin{table}[!htbp]
\centering\small
\caption{Positioning against representative defense approaches. The comparison describes problem scope, not an empirical performance ranking.}
\label{tab:related_comparison}
\begin{tabularx}{\textwidth}{p{0.24\textwidth}XXX}
\toprule
Approach & Intervention or emphasis & Relationship to this study & Comparison boundary\\
\midrule
Malware-oriented CPPS analysis \cite{xu2020analysis,xu2023resilience,zhang2025robustness} & Spreading, incubation, dependency, and protection rules & Supplies the malware and coupling motivation & Does not by itself validate this controller's state-continuation interface\\
Moving-target defense \cite{cho2020proactive,griffioen2021moving} & Changes system characteristics to impede adversarial inference or attacks & Demonstrates established adaptive cyber defense & Intervention space differs from selecting protected cyber nodes\\
Adaptive nonlinear power-system defense \cite{mohammadpourfard2025adaptive} & Adaptive moving-target defense in a dynamic power-system setting & Recent comparator for adaptive-defense positioning & No matched numerical benchmark is claimed here\\
Dynamic-botnet resiliency \cite{zhao2024iotbotnet} & Integrated cyber--physical resiliency under dynamic botnet attacks & Related cross-layer decision setting & Attack model and action space require alignment before numerical comparison\\
Present architecture & Rollout selection among interpretable protection sets with a baseline acceptance gate & Separates predicted gain, realized loss, and implementation contract & Service proxies and limited rolling validation; no global safety or superiority claim\\
\bottomrule
\end{tabularx}
\end{table}
